\chapter{Introduction}

\section{Research Background}

The banking system performs important functions in modern financial
markets, including capital allocation, credit creation, and liquidity
provision. Banks are not only connected with firms, households,
governments, and other sectors, but also conduct short-term funding,
liquidity adjustment, and risk sharing through the interbank market. The
banking system should not be regarded as a collection of independent
financial institutions. Instead, it can be understood as a complex
financial network formed by multiple creditor-debtor
relationships {[}1{]}.

Under normal market conditions, interbank transactions help improve the
efficiency of fund allocation and relieve short-term liquidity pressure
on individual institutions. However, when the market is affected by
external shocks or when some banks experience liquidity or solvency
problems, the interbank connections originally used for liquidity
sharing may also become channels for risk transmission {[}3{]}. In this
sense, interbank networks have a dual role: they can improve the
allocation of liquidity in stable periods, but they may also amplify
financial contagion during periods of stress.

Systemic risk is one of the core issues in financial stability research.
Unlike the individual risk of a single financial institution, systemic
risk focuses on how local shocks are transmitted through balance sheet
linkages, payment obligations, and market feedback among financial
institutions, and how these shocks further affect the stability of the
entire financial system. In an interbank market, this process is closely
related to creditor-debtor relationships. When a bank cannot repay its
interbank debt on time because of insufficient liquidity or capital, the
asset quality and cash flow of its creditor banks may also be affected.
If these creditor banks have weak capital buffers or insufficient
liquidity reserves, the initial shock may continue to spread to other
banks and eventually lead to chain defaults or widespread capital
stress.

The clearing process in a multilateral interbank debt network provides
an important mechanism for understanding this type of risk propagation.
In such a network, the final payment result is not simply determined by
whether one bank defaults individually. Instead, each bank's actual
payment ability depends not only on its own available resources, but
also on the payments it receives from other banks. The Eisenberg--Noe
clearing framework provides a theoretical basis for analyzing payment
relationships and default propagation in financial networks with
interdependent obligations {[}2{]}.

The 2008 global financial crisis further showed that complex debt
relationships and liquidity dependence among financial institutions can
significantly amplify initial shocks. After the collapse of Lehman
Brothers, confidence in financial markets declined rapidly, interbank
funding conditions deteriorated, market liquidity contracted sharply,
and risks spread from individual institutions to the broader financial
system. This experience indicates that analyzing the financial condition
of a single bank is not sufficient for assessing overall financial
stability. It is also necessary to consider interbank network
connections, debt clearing relationships, and market transaction
mechanisms {[}6{]}. More recent banking turmoil in 2023 also highlighted
the importance of liquidity risk, confidence shocks, and rapid stress
transmission in modern banking systems {[}25{]}.

In banking supervision, indicators such as the capital adequacy ratio,
liquidity coverage ratio, and net stable funding ratio are widely used
to measure bank soundness. The capital adequacy ratio reflects a bank's
ability to absorb losses, the liquidity coverage ratio reflects its
ability to cope with short-term cash flow pressure, and the net stable
funding ratio measures the stability of its medium- and long-term
funding sources. The Basel III framework strengthened capital and
liquidity regulatory requirements after the global financial
crisis {[}12--14{]}. However, in systemic risk analysis,
changes in individual bank regulatory indicators cannot fully describe
the process of risk propagation through a network. Whether a bank
contributes to systemic risk depends not only on its own balance sheet,
but also on its position in the interbank network, its creditor-debtor
relationships with other banks, and the mechanism through which funds
are reallocated in the market.

Interbank lending relationships add another difficulty: they are not
fixed. To simplify the analysis, many existing studies assume that the
interbank network structure is exogenously given, or analyze shock
propagation on a static network. However, in real markets, interbank
lending relationships are continuously formed and adjusted during the
trading process. These relationships are affected by liquidity
conditions, funding needs, interest rate requirements, risk preferences,
and counterparty selection {[}10, 11{]}. It is important to study how
interbank networks are formed under different loan matching mechanisms
and how this network formation process affects the propagation of
systemic risk.

The central question is whether alternative transaction designs change
how interbank exposures are formed and how quickly systemic stress is
realized when otherwise comparable banks face the same market
environment. The simulations keep bank types, initial balance-sheet
rules, regulatory thresholds, contract accounting, and the clearing
method common while changing the organization of matching. The
centralized benchmark performs market-wide rate-based allocation on each
business day, whereas decentralized RFQ operates through limited bilateral
enquiries and a locally observed counterparty pool. A system of 30
institutions---one central bank, 20 commercial banks, and nine shadow
banks---is followed in business-day steps, so a loan agreed today can
affect liquidity, repayment pressure, and clearing outcomes on later
days.

\section{Research Motivation}

Systemic risk in the interbank market emerges from the interaction of
bank-specific conditions, bilateral exposures, market liquidity, and
trading arrangements. A substantial part of the literature examines
contagion on a given network, focusing on how the default of one
institution affects others through existing claims {[}4{]}. This approach
provides the basis for analyzing clearing and default propagation, but
it treats the network itself as fixed. In practice, interbank exposures
change as banks respond to funding needs, liquidity constraints,
interest-rate conditions, and counterparty risk {[}10{]}. A static
network therefore omits the lending decisions through which exposures
are created and concentrated before clearing begins.

This study treats network formation as part of the risk process rather
than as an external input. Financial stress may accumulate during loan
matching, even before any bank defaults. Repeated borrowing from a
limited group of lenders can concentrate exposures, while unsuccessful
funding requests can leave some borrowers under liquidity pressure. The
identity of the matched counterparties, the amount of unmet demand, and
the distribution of lending across institutions may all influence which
banks become vulnerable later. The analysis therefore links the
formation of interbank loans with subsequent balance-sheet changes and
clearing outcomes.

The second motivation of the study is to compare how different interbank
transaction designs affect systemic risk. Centralized matching is used
as a high-coordination benchmark in which current orders are evaluated
jointly across the market on every business day. By contrast,
decentralized request-for-quote matching restricts borrowers to local
enquiries and bilateral quotations, which more closely reflects the fragmented and
over-the-counter nature of many interbank transactions {[}11{]}. Under
RFQ, system-wide liquidity may be sufficient even though an individual
borrower cannot obtain an acceptable quote. The comparison therefore
captures differences in information scope, counterparty access,
funding-completion rules, and pricing while keeping daily transaction timing and
the subsequent contract-accounting and clearing procedure common.

Liquidity pressure, capital stress, and payment contagion are examined
within a common framework, although they represent distinct forms of
financial distress. A bank may remain solvent while being unable to meet
a short-term payment, or it may maintain sufficient liquidity while
having only a limited capital buffer. Losses transmitted through
interbank claims can subsequently weaken both positions {[}15, 16{]}.
For this reason, the simulation tracks liquidity conditions, capital
adequacy, interbank assets and liabilities, contractual payments, and
clearing results. Systemic risk is then evaluated through a composite
indicator that includes the bank failure rate, the share of banks under
capital stress, and the aggregate capital gap.

The results may depend on how regulatory stress and aggregate risk are
measured. Raising the capital adequacy threshold applies a stricter
criterion and may classify a larger number of banks as
capital-constrained {[}14{]}. Similarly, changing the weights assigned to
failures, capital stress, and capital shortfalls alters the
responsiveness of the composite risk measure. The study therefore
conducts sensitivity tests for both the capital adequacy threshold and
the systemic-risk weights. Overall, the proposed framework compares
centralized matching and decentralized RFQ under otherwise consistent
bank characteristics, market conditions, contract accounting, and
clearing procedures. This design makes it possible to identify whether
differences in systemic-risk dynamics originate from the way interbank
exposures are formed.

\section{Research Purpose}

This study aims to develop a dynamic interbank network simulation
framework for examining whether different transaction designs lead to
different interbank network structures, funding outcomes, and
systemic-risk dynamics. In contrast to models that impose a fixed
interbank network, bilateral lending relationships are generated
endogenously from banks' changing balance-sheet conditions and trading
decisions.

The study compares a daily centralized market-wide allocation mechanism
with a daily decentralized request-for-quote mechanism under common bank
characteristics, market conditions, contract-accounting rules, and
due-payment clearing. Rollover and policy support are then switched on
and off to examine how repayment timing and intervention affect the
speed at which systemic stress emerges. Formal experiments use a
pre-specified 1,000-business-day horizon and stop early only if the number
of surviving non-central banks falls to one or fewer. Runs that do not
collapse by the horizon are treated as right-censored in survival-related
statistics.

The study addresses the following research questions:

\begin{enumerate}
\def\labelenumi{\arabic{enumi}.}
\item
  Do daily centralized matching and daily decentralized RFQ produce
  different interbank network structures and funding outcomes?
\item
  Do the two transaction designs generate different paths and timing of
  payment shortfalls, bank failures, capital stress, and systemic risk
  when they are compared over the same elapsed time?
\item
  How do rollover and policy support alter the timing and relative
  evolution of systemic risk, and do the qualitative conclusions remain
  stable under changes in the CAR threshold and the weights of the
  composite systemic-risk indicator?
\end{enumerate}

\section{Research Content and Methods}

The study constructs a discrete-time simulation model in which
heterogeneous banks adjust their balance sheets, form interbank loans,
maintain outstanding contracts, settle scheduled obligations, and update
their risk conditions over multiple business days. Interbank lending
relationships are generated endogenously rather than imposed as a fixed
network.

Two transaction mechanisms are implemented within the same balance-sheet
and clearing framework. The centralized mechanism executes market-wide
rate-based allocation every business day and requires complete borrower
funding before tentative trades are committed. The decentralized RFQ
mechanism also operates every business day, but limits each borrower to a
locally constructed lender pool and retains acceptable partial trades.
The resulting comparison therefore isolates differences in information
scope, counterparty discovery, pricing, and allocation more closely than
a design with unequal trading frequencies.

The rollover switch selects the repayment schedule when a trade is
originated. When rollover is enabled, each new loan is recorded as a
5--20-period amortizing installment contract. When it is disabled, the
loan is a next-day single-payment contract. Policy support can provide
next-period liquidity assistance and same-day post-project
solvency/capital assistance. The four ON/OFF combinations of rollover and
policy support are compared under both transaction mechanisms.

Systemic risk is evaluated using the bank failure rate, capital-stress
share, and capital-gap ratio, both separately and through a composite
indicator. Monte Carlo replications account for stochastic initial
conditions, market changes, shocks, and RFQ enquiries. The formal design
uses 20 paired replications over (T=1000) business days. Early stopping is
reserved for systemic collapse, defined operationally as one or fewer
surviving non-central banks; temporary network stability does not terminate
a formal run.

Sensitivity analyses examine the CAR threshold and the weights assigned
to the systemic-risk components, while a structural robustness check
varies the single-loan cap.

\section{Thesis Structure}

Chapter~1 introduces the research background, motivation, purpose,
research questions, and methodological approach.

Chapter~2 reviews the literature on systemic risk, financial contagion,
interbank networks, clearing mechanisms, prudential regulation, and
interbank matching, and identifies the research gap addressed by this
study.

Chapter~3 presents the overall model framework, including bank types,
state variables, market conditions, regulatory indicators, and the
dynamic simulation process.

Chapter~4 explains the daily centralized market-wide mechanism, the daily
decentralized RFQ mechanism, and the origination-based rollover and coupon
payment rules.

Chapter~5 describes the Eisenberg--Noe clearing procedure, the
bank-failure definition, and the measurement of systemic risk.

Chapter~6 specifies the simulation settings, collapse stopping rule,
bank population, feature scenarios, paired Monte Carlo procedure,
common-horizon comparison rule, and sensitivity analyses.

Chapter~7 reports and discusses the common-horizon results for the two
matching mechanisms, the rollover/support scenarios, network formation,
and the parameter-sensitivity experiments.

Chapter~8 summarizes the main findings, contributions, limitations, and
possible directions for future research.
